\begin{pastexam}[label={px:ntu-110-comm-dsa-45}]{臺大電信所 110 學年度（B）第 45 題}
（複選；每題 2.5 分，每答錯一個選項倒扣 0.5 分至該題 0 分止）

Consider a min-heap. Which of the following statements are true?
\begin{enumerate}[label=\Alph*.]
\item For any subtree of a min-heap, the root of the subtree contains the smallest key occurring anywhere in that subtree.
\item An array that is in ascending order is a binary min-heap.
\item The sequence [1,5,6,7,14,17,10,23,13,12] is a binary min-heap.
\item The largest key must reside in a leaf.
\item The smallest key must reside in a leaf
\end{enumerate}
\end{pastexam}

\begin{pastsoln}
答案是 A、B、D。

最小堆積的性質：每個節點的鍵 $\le$ 它子節點的鍵（資結書第~8.1~節）。 陣列從索引 $1$ 算，節點 $i$ 的子節點是
$2i$、$2i + 1$。
\begin{itemize}
\item A 對：沿著任何一條往下的路徑，鍵值不會變小（每一步都 $\ge$ 父節點），子樹裡任一節點都能從子樹的根
往下走到，所以都 $\ge$ 子樹的根。
\item B 對：由小到大排好的陣列，$A[i] \le A[2i]$、$A[i] \le A[2i + 1]$ 都成立（索引較小的值較小），而陣列本身
就是完整二元樹的存法，所以是最小堆積。
\item C 錯：逐一檢查父子：$A[5] = 14$ 的子節點是 $A[10] = 12$，$14 > 12$，違反堆積性質。 （其餘 $1 \le 5, 6$；
$5 \le 7, 14$；$6 \le 17, 10$；$7 \le 23, 13$ 都對。）
\item D 對：若最大的鍵在某個有子節點的位置，它的子節點的鍵 $\ge$ 它，也就是同樣是最大值；一路往下總會到葉子，
所以葉子裡一定有最大的鍵（鍵互不相同時，最大的鍵就只能在葉子）。
\item E 錯：最小的鍵在根。 只有一個節點時根同時是葉子，但一般情形不成立。
\end{itemize}
\end{pastsoln}
